\section{Introduction}\label{sec:intro}

Automatic Chord Recognition (ACR) underpins computational musicology, structural segmentation, and content-based audio indexing~\cite{pauwels2019twenty}. Although ACR architectures have progressed from handcrafted chroma templates and hidden Markov models to convolutional-recurrent networks and self-attention Transformers~\cite{korzeniowski2016feature,park2019btc,gulati2020conformer}, performance remains bounded by an empirical triadic plateau: classifiers reach $82\%$--$84\%$ accuracy on major and minor triads but degrade sharply on four-note tetrads and extended harmonies~\cite{pauwels2019twenty} (e.g., $77.6\%$ MajMin versus $58.8\%$ Tetrads weighted recall for an HMM on the canonical $1{,}217$-song benchmark~\cite{humphrey2015four}).

This ceiling stems not from insufficient model capacity, but from the severe Zipfian skew of real-world musical corpora: in the $1{,}217$-song benchmark, major ($65.75\%$) and minor ($22.16\%$) triads occupy $87.91\%$ of non-silent duration ($73.7\,\text{h}$, $11{,}436{,}357$ frames). Rare triads ({\small\texttt{sus4}}, {\small\texttt{dim}}, {\small\texttt{sus2}}, {\small\texttt{aug}}) collectively span under $3.5\%$ of duration, while altered upper extensions ({\small\texttt{\#11}} at $0.13\%$, {\small\texttt{b9}} at $0.13\%$, {\small\texttt{b13}} at $0.08\%$) inhabit a sparse tail. Moreover, human inter-annotator disagreement exceeds $40\%$ on inversions and upper extensions~\cite{koops2019annotator}, compounding sample scarcity with label noise.

Recent architectures address large-vocabulary recognition via structured multi-task losses~\cite{mcfee2017structured}, 6-head chord decomposition~\cite{jiang2019large,akram2026chordformer}, and contrastive pseudo-labeling~\cite{li2025large}. Specifically, Jiang et al.~\cite{jiang2019large} and Akram et al.~\cite{akram2026chordformer} factorize chords into six categorical heads (root+triad, bass, and $7\text{th}, 9\text{th}, 11\text{th}, 13\text{th}$ degrees) trained under inverse-frequency reweighted cross-entropy. Yet loss reweighting cannot overcome an absence of physical acoustic exemplars: when an altered degree occurs in only a few dozen frames across a corpus, scaling its gradient penalty magnifies gradient variance and overfits to the idiosyncratic instrumentation and room acoustics of those few tracks. Resolving the long tail requires intervening directly on the training audio distribution.

We propose a Sim-to-Real pretraining framework driven by procedurally synthesized multitrack audio, contributing: (1) \textbf{Probabilistic Composition Generator}, an open-source engine sampling progressions along empirical transitions curated via the ChoCo aggregator~\cite{deberardinis2023choco} under deterministic quota scheduling (upweighting rare triads up to $15.6\times$ and altered extensions up to $12.3\times$), enforcing physical register bounds and $100\%$ extension sounding across four stems with random-mix augmentation to render $5{,}000$ tracks ($460.3\,\text{h}$) and frame-exact Harte labels~\cite{harte2005symbolic}; (2) \textbf{Two-Stage Sim-to-Real Curriculum}, a transfer regimen for ChordFormer~\cite{akram2026chordformer} combining synthetic pretraining with real-audio finetuning via inverted discriminative learning rates ($\mathbf{s}_{\text{lr}} = [0.2, 0.5, 1.0]$) and floor curriculum replay ($\rho(\eta) = 0.150 \!\to\! 0.100$); and (3) \textbf{Multi-Head and WCSR Benchmark Analysis}, a 5-fold cross-validated evaluation over all $11{,}436{,}357$ frames demonstrating best-in-class rare-chord recall ($+6.64\,\text{pp}$ on rare triads, $+25.16\,\text{pp}$ on natural 9ths, $79\times$ on {\small\texttt{\#9}}, $7.8\times$ on {\small\texttt{b13}}) alongside an analysis of precision-recall trade-offs, duration-weighted WCSR bias, and polyphonic overtone masking.
